\section*{Chapitre \ref{ch:qm:point-processes-and-hawkes-processes} --- Processus ponctuels et processus de Hawkes}

\begin{solution}{exo:qm:point-processes-and-hawkes-processes:1}
$n = 0.7$~; $\bar\lambda = 0.1/0.3 = 0.333$ par seconde~; $0.333 \times 23\,400 = 7\,800$ événements.
\end{solution}

\begin{solution}{exo:qm:point-processes-and-hawkes-processes:2}
$1/(1 - 0.7) = 3.33$ événements par immigrant~; les immigrants représentent $1 - n = 30\%$ des événements.
\end{solution}

\begin{solution}{exo:qm:point-processes-and-hawkes-processes:3}
$\Lambda(t) = at + bt^2/2$~; $\P(N_T = 0) = e^{-aT - bT^2/2}$.
\end{solution}

\begin{solution}{exo:qm:point-processes-and-hawkes-processes:4}
$\ln 2/\beta = 0.69$ secondes~; le saut $\alpha = 0.7$ vaut sept fois l'intensité de base $0.1$.
\end{solution}

\begin{solution}{exo:qm:point-processes-and-hawkes-processes:5}
Avec $\gamma = \beta(1 - n) = 0.3$ et $n(2 - n)\beta/(2(1 - n)) = 1.517$~: $1 + 2 \times 1.517\,(1/0.3 - (1 - e^{-0.3})/0.09) = 2.38$ à
une seconde, 10.5 à une minute, et $1/0.3^2 = 11.1$ à la limite.
\end{solution}

\begin{solution}{exo:qm:point-processes-and-hawkes-processes:6}
Moyenne 2, variance $\E[\Lambda] + \Var(\Lambda) = 2 + 1 = 3$. C'est un \omterm{def:qm:jump-processes:poisson}{processus de Poisson} mélangé (\omterm{def:qm:point-processes-and-hawkes-processes:cox}{processus de Cox})~: il est surdispersé sans aucune
auto-excitation, ce qui explique que la surdispersion seule ne permette pas d'identifier un \omterm{def:qm:point-processes-and-hawkes-processes:hawkes}{processus de Hawkes}.
\end{solution}

\begin{solution}{exo:qm:point-processes-and-hawkes-processes:7}
Moyenne 0.700, écart-type 0.011 sur quarante jours d'environ 7\,800 événements chacun.
\end{solution}

\begin{solution}{exo:qm:point-processes-and-hawkes-processes:8}
Une intensité de base constante ne peut pas suivre le profil intrajournalier de l'activité, et l'ajustement explique les périodes d'affluence par
une auto-excitation à noyau long~: une journée sans aucune auto-excitation mais avec un taux en forme de U donne 0.87. Modéliser la
variation de l'intensité de base (ou rééchelonner le temps selon le profil intrajournalier) avant d'interpréter le \omterm{def:qm:point-processes-and-hawkes-processes:hawkes}{ratio de branchement}, et vérifier
les résidus selon le moment de la journée.
\end{solution}

\begin{solution}{pb:qm:point-processes-and-hawkes-processes:1}
\textbf{1.} $\alpha/\beta = 0.7$.
\textbf{2.} $0.333$ par seconde~; 7\,800 transactions.
\textbf{3.} $3.33$~; 30~\% d'immigrants.
\textbf{4.} $\ln 2 = 0.69$ secondes.
\textbf{5.} 2.38 et 10.5.
\textbf{6.} 8\,006.
\textbf{7.} $\hat\mu = 0.103$, $\hat\alpha = 0.696$, $\hat\beta = 0.995$~: $\hat n = 0.700$.
\textbf{8.} $0.700 \pm 0.011$.
\textbf{9.} Ajustement de Hawkes~: moyenne 1.000, variance 0.98. Ajustement de Poisson~: moyenne égale à un par construction, variance 4.63.
\textbf{10.} 69.9~\% (7\,498 événements).
\textbf{11.} Trois fois plus élevé à l'ouverture et à la clôture~; taux moyen $0.333$ par seconde, $0.2$ à midi.
\textbf{12.} 7\,730.
\textbf{13.} $\hat n = 0.87$ avec $\hat\beta = 0.0076$ par seconde~: un noyau d'une durée d'environ deux minutes.
\textbf{14.} Moyenne 0.86 sur vingt jours saisonniers~; 0.04 sur vingt jours de Poisson à taux constant.
\textbf{15.} Un taux variant lentement fait se regrouper les événements dans le temps~; avec une intensité de base constante, la seule façon dont le modèle
peut placer davantage d'événements là où ils sont plus nombreux passe par une auto-excitation à noyau long.
\textbf{16.} Estimer une intensité de base variable dans le temps $\mu(t)$ conjointement avec le noyau, ou rééchelonner le temps selon le profil intrajournalier estimé
au préalable~; vérifier séparément les résidus à l'ouverture, à midi et à la clôture.
\textbf{17.} Que le \omterm{def:qm:point-processes-and-hawkes-processes:hawkes}{ratio de branchement} dépend de la forme du noyau, de la fenêtre d'estimation, de l'intensité de base et des valeurs aberrantes~:
c'est un nombre fragile, et les comparaisons d'une période à l'autre exigent la même méthode.
\textbf{18.} Prévoir l'activité à court terme et le flux toxique (élargir ou retirer les cotations lorsque l'intensité monte brusquement),
pour programmer l'exécution à l'écart des rafales, et pour générer des flux d'ordres réalistes pour les backtests.
\textbf{19.} Résultat nommé~: \emph{la bande auto-excitée}~: sur les jours à \omterm{def:qm:point-processes-and-hawkes-processes:hawkes}{ratio de branchement} vrai de 0.7, l'ajustement
retrouve $0.700 \pm 0.011$~; sur les jours sans excitation mais à intensité de base en forme de U, il annonce 0.87 (0.86 en moyenne),
contre 0.04 pour les jours de Poisson à taux constant.
\textbf{20.} Seulement lorsque les variations propres de l'intensité de base sont modélisées et que les résidus rééchelonnés en temps passent les contrôles.
\end{solution}

\begin{solution}{iq:qm:point-processes-and-hawkes-processes:1}
Un \omterm{def:qm:point-processes-and-hawkes-processes:pp}{processus de comptage} dont l'intensité vaut une intensité de base plus un noyau décroissant sommé sur les événements passés, de sorte que chaque événement
relève le taux des événements suivants. Le \omterm{def:qm:point-processes-and-hawkes-processes:hawkes}{ratio de branchement}, intégrale du noyau, est le nombre moyen d'événements
que chaque événement en déclenche, et la part à long terme de l'activité qui est endogène.
\iqlookfor{la formule de l'intensité et les deux lectures du \omterm{def:qm:point-processes-and-hawkes-processes:hawkes}{ratio de branchement}.}
\end{solution}

\begin{solution}{iq:qm:point-processes-and-hawkes-processes:2}
Prenons les espérances en régime stationnaire~: $\bar\lambda = \mu + n\bar\lambda$, donc $\bar\lambda = \mu/(1 - n)$. Pour $n \ge 1$, chaque événement engendre en
moyenne au moins un de plus, les grappes sont infinies avec une probabilité positive, et aucune version stationnaire
n'existe~: le processus explose.
\iqlookfor{l'argument du point fixe et la criticité.}
\end{solution}

\begin{solution}{iq:qm:point-processes-and-hawkes-processes:3}
Par \omterm{def:qm:point-processes-and-hawkes-processes:thinning}{amincissement}~: l'intensité exponentielle décroît entre deux événements, si bien que l'intensité courante la majore jusqu'au
prochain événement~; proposer à ce taux et accepter avec probabilité égale au rapport de l'intensité au majorant. Ou générer les immigrants,
puis des nombres de Poisson de descendants, génération par génération.
\iqlookfor{l'\omterm{def:qm:point-processes-and-hawkes-processes:thinning}{amincissement} d'Ogata ou la représentation en grappes, et pourquoi ces deux approches sont exactes.}
\end{solution}

\begin{solution}{iq:qm:point-processes-and-hawkes-processes:4}
L'intensité en $t_i$ requiert $A_i = \sum_{j<i}e^{-\beta(t_i - t_j)}$, et $A_i = e^{-\beta(t_i - t_{i-1})}(1 + A_{i-1})$~; l'intégrale
de l'intensité est une somme d'un terme par événement. Les deux se calculent en un seul parcours des données.
\iqlookfor{la récurrence.}
\end{solution}

\begin{solution}{iq:qm:point-processes-and-hawkes-processes:5}
Comme prévision à court terme du niveau d'activité des prochaines secondes et de son côté~: élargir ou décaler les cotations lorsque
l'intensité du côté opposé monte brusquement, car les rafales portent une sélection adverse~; dimensionner les limites d'inventaire selon le
flux attendu~; et, dans le simulateur, générer un flux groupé au lieu d'un flux de Poisson.
\iqlookfor{l'intensité comme prévision du flux et de la sélection adverse.}
\end{solution}

\begin{solution}{iq:qm:point-processes-and-hawkes-processes:6}
Rééchelonner le temps par le \omterm{def:qm:point-processes-and-hawkes-processes:intensity}{compensateur} ajusté et tester que les intervalles sont des exponentielles unitaires indépendantes~: moyenne et
variance, un graphique QQ, un test de Kolmogorov--Smirnov (chapitre~12), l'autocorrélation des intervalles, et le tout selon le moment
de la journée et hors échantillon.
\iqlookfor{le théorème de rééchelonnement du temps comme outil d'adéquation.}
\end{solution}
